% Part: methods
% Chapter: proofs
% Section: introduction

\documentclass[../../../include/open-logic-section]{subfiles}

\begin{document}

\olfileid{mth}{prf}{int}

\olsection{引言}

根据你在入门逻辑课程中的经验，你可能已经熟悉某种推演系统——也许是自然演绎或 Fitch 风格的推演系统，或者是证明树系统。你可能还记得在这些系统中做证明，要么证明一个公式，要么证明一个给定的论证是有效的。为此，你应用系统的规则，直到得到想要的结果。在对逻辑本身进行推理时，我们同样要证明一些事情，但在大多数情况下，我们并不使用推演系统。事实上，我们考虑的大多数证明都是用英语写成的（或许会夹杂一些符号语言），而非完全用一阶逻辑的语言。在构造此类证明时，起初你可能会感到无从下手——没有推演系统，我该如何证明某个命题？如何开始？如何知道我的证明是否正确？

在尝试证明之前，重要的是先知道证明是什么以及如何构造证明。顾名思义，一个\emph{证明}旨在表明某件事为真。你可以把它想象成一场对话——有人问你某件事是否为真，例如，除二以外的所有素数是否都是奇数。仅仅回答“是”是不够的；他们可能想知道\emph{为什么}。此时，你就会给他们一个证明。

在日常话语中，或许含糊地给出一个答案或不完整的答案就足够了。然而，在逻辑和数学中，我们需要严格的证明——我们想要排除\emph{任何}怀疑地表明某件事为真。这意味着我们证明中的每一步都必须有依据，而且该依据必须具有说服力（即，你所使用的假设确实是你所证定理陈述中假定的，你所应用的定义必须被正确地应用，你所援引的依据必须是正确的推理，等等）。

通常，我们证明的是某个陈述。我们用不同的名称来称呼我们要证明的陈述：命题、定理、引理或推论。命题是一个基本的、值得证明的陈述：重要到值得记录，但或许并非特别深刻或应用广泛。定理是一个重要的、意义重大的命题。其证明常分为若干步骤，有时以首次证明它的人命名（例如，Cantor（康托尔）定理、Löwenheim–Skolem（勒文海姆–斯科伦）定理）或以相关事实命名（例如，完备性定理）。引理是用于证明更重要结果的命题或定理。有些令人困惑的是，有时引理本身也是重要的结果，并且也以引入它们的人命名（例如，Zorn（佐恩）引理）。推论是容易从另一结果推出的结果。

待证的陈述通常包含一些假设，用以明确我们正在对哪类事物进行证明。它可能以“设 $!A$ 是一个形如 $!B \lif !C$ 的公式”或“假设 $\Gamma \Proves !A$”之类的形式开始。这些是命题、定理或引理的\emph{假设}，在你的证明中，你可以假定这些为真。它们限定了我们正在证明的内容，同时也为我们讨论的对象引入了一些名称。例如，如果你的命题以“设 $!A$ 是一个形如 $!B \lif !C$ 的公式”开头，那么你证明的仅仅是关于某类公式（即条件句）的某个性质，并且 $!B \lif !C$ 被理解为一个任意的条件句，你的证明将围绕它展开。

\end{document}